% GENERATED FILE. Edit canonical lessons, then run npm run generate:books.
% canonical-source-hash: fnv1a64:e9c420c4e46128da
% canonical-lessons: ML-C23-divasam, ML-C23-naal

\chapter{Day}
\label{ch:ml-day}

\begin{chapteropening}
\textbf{By the end of this chapter:} \emph{I can use divasam for an everyday day in Malayalam and still recognise native nāḷ where it survives.}

Find nāḷ alive inside nāḷe, oru nāḷ and piṟannāḷ, and answer whether Sanskrit divasam erased it.
\end{chapteropening}

\section[\ml{ദിവസം}]{\ml{ദിവസം} (divasam) — the everyday Sanskrit day-word}

\label{lesson:ML-C23-divasam}



\begin{quote}
You already met Malayalam's native word for \textquotedblleft{}day\textquotedblright{} back in Chapter 4 — hiding inside \textquotedblleft{}tomorrow.\textquotedblright{} This lesson gives it, and its Sanskrit rival, their own honest reckoning.
\end{quote}

\subsection*{What to know first}
\textbf{\ml{ദിവസം}} (\textbf{divasam}) — \textquotedblleft{}\textbf{day}\textquotedblright{} — is Malayalam's common, everyday word for counting a day (as in \textbf{\ml{മൂന്ന്} \ml{ദിവസം}}, \emph{mūnnu divasam}, \textquotedblleft{}three days\textquotedblright{}). It's a Sanskrit \textbf{tatsama}, borrowed whole from \textbf{\dv{दिवस}} (\emph{divasa}), which itself traces to \textbf{Proto-Indo-European} \emph{\textbf{*dyew-}} (\textquotedblleft{}\textbf{to shine}\textquotedblright{}), via the Sanskrit root \textbf{\dv{दिव्}} (\emph{div}). That's the \textbf{same ultimate PIE root} behind Hindi/Kannada's \textbf{\dv{दिन}/\ka{ದಿನ}} (\emph{dina}) and Latin's own \emph{diēs} — but be careful here: \emph{divasa} and \emph{dina} are \textbf{two different Sanskrit formations} off that shared root (one built with an \emph{-asa} suffix, one with Sanskrit's own \emph{-nó-} formation), not the same word wearing two scripts. Same root, different branch — the \emph{din}/\emph{diēs} pattern, replaying itself one level deeper, inside Sanskrit.

\begin{cousinweb}
Malayalam also has \textbf{\ml{ദിനം}} (\textbf{dinam}) — the \textbf{exact same} Sanskrit tatsama word as Kannada's plain, everyday \emph{dina}. But in Malayalam, \textbf{dinam} is the more \textbf{literary, formal} alternate (think \textbf{\ml{ജന്മദിനം}}, \emph{janmadinam}, \textquotedblleft{}birthday,\textquotedblright{} the modern calque-style term) — \textbf{divasam} does the everyday counting work instead.
\end{cousinweb}

\subsection*{Your turn}
Say these aloud:

\begin{itemize}
\raggedright
  \item \textquotedblleft{}divasam\textquotedblright{} — \textquotedblleft{}day,\textquotedblright{} the everyday counting word
  \item \textquotedblleft{}dinam\textquotedblright{} — the same Sanskrit word as Kannada's dina, but formal here
\end{itemize}

\begin{tcolorbox}[breakable,colback=teal!4,colframe=teal!35!black,title={Before you move on}]
What PIE root does \textbf{\ml{ദിവസം}} ultimately share with Hindi/Kannada's \emph{dina} and Latin's \emph{diēs}? (\emph{\textbf{*dyew-}}, \textquotedblleft{}to shine\textquotedblright{} — though \emph{divasa} and \emph{dina} are different Sanskrit formations off it, not the same word.) Which is the everyday counting word? (\textbf{divasam}.) Which is more literary or formal? (\textbf{dinam}.)
\end{tcolorbox}
\section[\ml{നാൾ}]{\ml{നാൾ} — the native day-word that narrowed its job}

\label{lesson:ML-C23-naal}



\begin{quote}
\textbf{\ml{ദിവസം}} counts ordinary days and \textbf{\ml{ദിനം}} writes formally. But Malayalam’s native day-word never disappeared—you met it long ago inside “tomorrow.”
\end{quote}

\subsection*{What to know first}
\textbf{\ml{നാൾ}} (\emph{nāḷ}) means “day” and descends from Proto-Dravidian \emph{\textbf{*nāḷ}}, cognate with Tamil \textbf{\ta{நாள்}} and Kannada \textbf{\ka{ನಾಳು}}. Chapter 4 already gave you \textbf{\ml{നാളെ}} (\emph{nāḷe}, “tomorrow”), with the native day-word still visible inside.

The Sanskrit loan \textbf{\ml{ദിവസം}} took over ordinary counting, but \textbf{\ml{നാൾ}} narrowed rather than vanished. \textbf{\ml{ഒരു} \ml{നാൾ}} (\emph{oru nāḷ}, “one day”) remains natural in storytelling, and the word stays productive in \textbf{\ml{പിറന്നാൾ}} (\emph{piṟannāḷ}), “birthday,” literally “day of birth.” That traditional compound stands beside the newer Sanskrit-style \textbf{\ml{ജന്മദിനം}} (\emph{janmadinam}).

Three words therefore carry three honest jobs: \emph{divasam} counts, \emph{dinam} is formal, and native \emph{nāḷ} tells stories and marks birthdays. This echoes Kannada’s native \emph{hagalu}: borrowing narrowed a native word’s domain without erasing it.

\subsection*{Your turn}
Say these aloud:

\begin{itemize}
\raggedright
  \item \emph{nāḷe} — tomorrow, with \emph{nāḷ} inside
  \item \emph{oru nāḷ} — one day, narrative register
  \item \emph{piṟannāḷ} — birthday, day of birth
\end{itemize}

\begin{tcolorbox}[breakable,colback=teal!4,colframe=teal!35!black,title={Before you move on}]
Has Sanskrit \textbf{divasam} erased native \textbf{\ml{നാൾ}}? (\textbf{No.}) Where does \emph{nāḷ} live now? (\textbf{Inside tomorrow, in narrative “one day,” and in the living compound for birthday.})
\end{tcolorbox}
